% !TeX program = pdflatex
\documentclass[11pt,a4paper]{amsart}

\newcommand{\notelanguage}{finnish}
% Shared preamble for course notes and exercise sets.

\usepackage[T1]{fontenc}
\usepackage{lmodern}
\usepackage[english]{babel}

\usepackage[a4paper,margin=30mm]{geometry}
\usepackage{microtype}
\usepackage{graphicx}
\usepackage{enumitem}

\usepackage{amsmath}
\usepackage{amssymb}
\usepackage{amsthm}
\usepackage{mathtools}
\usepackage{aliascnt}

\usepackage[hidelinks,pdfusetitle]{hyperref}

\numberwithin{equation}{section}
\setlist{itemsep=0.25em,topsep=0.5em}

\theoremstyle{plain}
\newtheorem{theorem}{Theorem}[section]
\newaliascnt{lemma}{theorem}
\newtheorem{lemma}[lemma]{Lemma}
\aliascntresetthe{lemma}
\newaliascnt{proposition}{theorem}
\newtheorem{proposition}[proposition]{Proposition}
\aliascntresetthe{proposition}
\newaliascnt{corollary}{theorem}
\newtheorem{corollary}[corollary]{Corollary}
\aliascntresetthe{corollary}

\theoremstyle{definition}
\newaliascnt{definition}{theorem}
\newtheorem{definition}[definition]{Definition}
\aliascntresetthe{definition}
\newaliascnt{example}{theorem}
\newtheorem{example}[example]{Example}
\aliascntresetthe{example}
\newaliascnt{problem}{theorem}
\newtheorem{problem}[problem]{Problem}
\aliascntresetthe{problem}

\theoremstyle{remark}
\newaliascnt{remark}{theorem}
\newtheorem{remark}[remark]{Remark}
\aliascntresetthe{remark}

\usepackage[nameinlink,noabbrev]{cleveref}

\crefname{theorem}{theorem}{theorems}
\crefname{lemma}{lemma}{lemmas}
\crefname{proposition}{proposition}{propositions}
\crefname{corollary}{corollary}{corollaries}
\crefname{definition}{definition}{definitions}
\crefname{example}{example}{examples}
\crefname{problem}{problem}{problems}
\crefname{remark}{remark}{remarks}

\DeclareMathOperator{\im}{im}
\DeclareMathOperator{\rank}{rank}
\DeclarePairedDelimiter{\abs}{\lvert}{\rvert}
\DeclarePairedDelimiter{\norm}{\lVert}{\rVert}
\DeclarePairedDelimiter{\set}{\lbrace}{\rbrace}


\usepackage{tikz}

\title{Analyysi I\\Ryhmätyö 3 ratkaisut}
\author{}
\date{}

\begin{document}

\exerciseheading{Analyysi I}{Ryhmätyö 3 ratkaisut}

\begin{exercise}
  \leavevmode\par
\begin{enumerate}[label=\alph*), nosep]
  \item Määritä seuraavien joukkojen supremum ja infimum:
    \begin{itemize}[nosep]
      \item \([1,7]\) ja \((1,7)\).
      \item \(\left\{\frac{2}{n + 1} \colon n = 1,2,3,\ldots\right\}\).
      \item \(\left\{\frac{3n}{n - 2} \colon n = 3,4,5,\ldots\right\}\).
    \end{itemize}
  \item Onko edellisissä joukoissa suurinta tai pienintä alkiota?
  \item Selitä omin sanoin, miten seuraavat käsitteet eroavat toisistaan:
    supremum, yläraja, suurin alkio. Esitä käsitteet myös graafisesti
    tehtävän (a) viimeiselle joukolle.
\end{enumerate}
\end{exercise}

\begin{solution}
Merkitään \(A = \set{\frac{2}{n + 1}\colon n \in \NNone}\) ja
\(B = \set{\frac{3n}{n - 2} \colon n = 3,4,5,\cdots}\).
\begin{enumerate}[label=\alph*), nosep]
\item
  \begin{itemize}[nosep]
    \item Joukon \([1, 7]\) suurin alkio on \(7\), joten \(\sup{[1, 7]} = 7\).
      Joukon \([1, 7]\) pienin alkio on \(1\), joten \(\inf{[1, 7]} = 1\).
      Joukossa \((1, 7)\) ei ole pienintä, eikä suurinta alkiota, mutta joukko
      sisältää lukuja mielivaltaisen lähellä lukuja \(1\) ja \(7\). Täten,
      \(\sup{(1, 7)} = 7\) ja \(\inf{(1, 7)} = 1\).
    \item Lukujono \((x_n)\), jossa \(x_n = \frac{2}{n + 1}\) on vähenevä ja
      suppenee kohti lukua \(0\). Täten
      \(\inf{A} = 0\). Koska \((x_n)\) on
      vähenevä ja \(x_1 = 1\), on jonon arvojoukon \(A\) suurin alkio \(1\),
      joka on tällöin myös \(\sup{A} = 1\).
    \item Lukujono \((x_n)\), jossa \(x_n = \frac{3n}{n - 2}\) on vähenevä ja
      suppenee kohti lukua \(3\). Täten \(\inf{B} = 3\). Koska \((x_n)\)
      on vähenevä ja \(x_3 = 9\), on jonon arvojoukon suurin alkio \(9\),
      eli \(\sup{B} = 9\).
  \end{itemize}
\item Joukolla \([1, 7]\) on suurin ja pienin alkio. Joukolla \((1, 7)\) ei ole
  suurinta, eikä pienintä alkiota. Oletetaan, että \(k\) on joukon
  \((1, 7)\) suurin alkio. Luku \(k + \frac{7 - k}{2}\) on kuitenkin lukua \(k\)
  suurempi, mutta kuitenkin pienempi kuin \(7\). Tämä on ristiriita.
  Todistus sille, ettei joukolla ole pienintä alkiota on samanlainen.
  Joukolla \(A\) on suurin alkio \(\max{A} = 1\), mutta ei pienintä alkiota.
  Joukolla \(B\) on suurin alkio \(\max{B} = 9\), mutta ei pienintä alkiota.
  Kummankin joukon jokaista alkiota pienempi alkio saadaan
  seuraavalla indeksillä:
  \begin{gather*}
    \frac{2}{n+2}<\frac{2}{n+1}\quad(n\ge1),\\[0.5em]
    \frac{3(n+1)}{n-1}=3+\frac{6}{n-1}
    <3+\frac{6}{n-2}=\frac{3n}{n-2}\quad(n\ge3).
  \end{gather*}
  Siksi mikään alkio ei voi olla joukkonsa pienin.
\item Yläraja on luku, jota yksikään joukon alkio ei ylitä.
  Supremum on pienin yläraja. Suurin alkio on joukkoon kuuluva yläraja,
  joten se on myös supremum. Supremumin ei kuitenkaan tarvitse kuulua joukkoon.
  Joukolle \(B\) pätee \(\sup B=\max B=9\), ja sen ylärajat ovat \([9,\infty)\).
  \begin{center}
    \begin{tikzpicture}[x=1.15cm,y=1cm]
      \draw[->] (2.6,0) -- (10.5,0);
      \foreach \x in {3,4,5,6,7,8,9,10} {
        \draw (\x,0.06) -- (\x,-0.06) node[below] {\small \x};
      }
      \foreach \n in {3,...,18} {
        \fill ({3+6/(\n-2)},0) circle (1.6pt);
      }
      \node[above] at (3.2,0) {\(\cdots\)};
      \draw[fill=white] (3,0) circle (2pt);
      \node[above] at (6,0.05) {\small Joukon \(B\) alkioita};
      \draw[thick,->] (9,0.8) -- (10.5,0.8);
      \fill (9,0.8) circle (2pt);
      \draw[dashed] (9,0.1) -- (9,0.7);
      \node[above] at (9.7,0.8) {\small Ylärajat};
      \node[below] at (9,-0.5) {\(\sup B=\max B=9\)};
    \end{tikzpicture}
  \end{center}
\end{enumerate}
\end{solution}

\begin{exercise}
Olkoot joukot \(B_n \subset (-\frac{2}{n}, \frac{3}{n})\) kaikilla
\(n \in \NNone\).
\begin{enumerate}[label=\alph*), nosep]
\item Onko lukujono \((\sup B_n)\) välttämättä kasvava tai vähenevä?
\item Osoita, että lukujono \((\sup B_n - \inf B_n)\) suppenee kohti nollaa.
\end{enumerate}
\end{exercise}

\begin{solution}
Oletetaan, että \(B_n \neq \varnothing\) kaikilla \(n \in \NNone\).
\begin{enumerate}[label=\alph*), nosep]
\item Ei. Olkoon \(B_1 = (0, 1)\), \(B_2 = \left(0, \frac{3}{2}\right)\) ja
\(B_n = \left(0, \frac{1}{n}\right)\) kun \(n \geq 3\). Tällöin
\(B_n \subset \left(-\frac{2}{n}, \frac{3}{n}\right)\) kaikilla \(n
\in \NNone\), mutta \(\sup{B_1} = 1 < \sup{B_2} = \frac{3}{2}\) ja
\(\sup{B_n} > \sup{B_{n + 1}}\) kaikilla \(n \geq 3\). Täten
\((\sup{B_n})\) ei ole kasvava, eikä vähenevä.
\item Olkoon \((x_n)\) lukujono, jolle \(x_n = \sup{B_n} - \inf{B_n}\). Koska
\(\sup{\left(-\frac{2}{n}, \frac{3}{n}\right)} = \frac{3}{n}\) ja
\(B_n \subset \left(-\frac{2}{n}, \frac{3}{n}\right)\), on oltava
\(\sup{B_n} \leq \frac{3}{n}\). Samalla päättelyllä saadaan
\(-\frac{2}{n} \leq \inf{B_n}\). Tiedetään myös \(\inf{B_n} \leq \sup{B_n}\).
Täten, edellisten arvioiden nojalla
\[
\abs{x_n} = \sup{B_n} - \inf{B_n} \leq \frac{3}{n} + \frac{2}{n} = \frac{5}{n}.
\]
Olkoon \(\epsilon > 0\). Valitaan \(N \in \NNone\) siten, että
\(N > \frac{5}{\epsilon}\). Tällöin,
\[
\abs{x_n} \leq \frac{5}{n} < \epsilon,
\]
kun \(n > N\), eli \(x_n \to 0\).
\end{enumerate}
\end{solution}

\end{document}
